\section{ELF 页面怎样从文件进入内存}

文件映射不是把 \code{read} 换一个名字，而是让文件字节成为页故障可以反复
取得的后备。VMA、PTE、后备对象与缓存页必须分别表达地址意图、当前驻留、
稳定来源和可丢弃内容。

\topic{从主动装入到 fault-driven execution}
传统 loader 可以在创建进程时遍历所有 ELF \code{PT\_LOAD}，为每页分配
frame、读取文件并建立 PTE。这条 eager 路径容易理解，却把启动时间和内存
占用绑定到整个映像大小。程序可能只执行入口附近几页，未触及的错误处理与
只读表格仍被提前装入。

虚拟内存允许把这项工作推迟：

\begin{enumerate}[label=\arabic*.]
  \item ELF reader 先完整校验 header、segment、范围、权限和入口；
  \item Kernel 建立文件后备与 \code{ExecutableImage} VMA；
  \item 返回用户态前只确保入口 instruction page 可执行；
  \item CPU 第一次访问其余页时产生用户 not-present \#PF；
  \item fault handler 由 VMA 找到文件 offset，读取或命中 cache；
  \item PTE 发布后返回，处理器重试原来的 load/store/fetch。
\end{enumerate}

校验不能延迟到每次 fault。若在进程已经运行后才发现 ELF segment 重叠或
入口无效，候选映像的原子提交就被破坏。这里只延迟物理装入，不延迟结构
合法性判断。

\topic{四类状态不能压进一个对象}\par
\begin{osgridlongtable}{L{0.20\textwidth}L{0.31\textwidth}L{0.33\textwidth}}
\textbf{对象} & \textbf{保存的事实} & \textbf{不应保存的事实}\\ \hline
VMA & 半开地址区间、R/W/X、kind、后备与文件区间 & 当前 frame 和硬件
accessed 位\\ \hline
PTE & present、frame、U/S、R/W、NX & fd、inode 生命周期与文件长度\\ \hline
FileBacking & 可按 offset 读取的稳定 VFS 打开实例或内存镜像 & 某页是否
已经驻留\\ \hline
CachePage & 文件页身份、clean frame、映射引用与 LRU 代数 & 虚拟区间
split 和进程权限\\ \hline
\end{osgridlongtable}

PTE 不存在正是 fault 发生的前提，因此来源不能只存于 PTE。fd 可以关闭和
复用，因此来源也不能是进程局部整数。CachePage 可能被多个 VMA 引用，不能
反过来拥有某一个 VMA 的地址几何。

\topic{稳定文件身份来自对象代数}
文件页键由

\begin{center}
\code{(superblock id, superblock generation, inode id, inode generation,
page index)}
\end{center}

构成。superblock 区分挂载实例，inode 区分文件，两个 generation 阻止同号
和同槽复用。page index 是文件 offset 除以 4096。fd 只在 \code{mmap}
建立时用于找到打开实例；FileBacking 取得独立引用后，用户可以立即
\code{close(fd)}，映射继续有效。

\begin{pdfdiagram}
  \node[os state] (fd) {Process FileTable\\fd};
  \node[os block,right=of fd] (description) {FileDescription\\VFS OpenFile};
  \node[os state,right=of description] (backing) {FileBacking\\owner + generation};
  \node[os block,below=of backing] (vma) {file VMA\\offset + length};
  \node[os state,left=of vma] (cache) {CachePage\\identity + frame};
  \draw[os arrow] (fd) -- node[above,font=\scriptsize]{lookup} (description);
  \draw[os arrow] (description) -- node[above,font=\scriptsize]{retain} (backing);
  \draw[os arrow] (backing) -- (vma);
  \draw[os arrow] (vma) -- node[below,font=\scriptsize]{fault} (cache);
\end{pdfdiagram}
\diagramnote{close 只删除 fd 引用；munmap、exec 或 exit 才释放 VMA 持有的
后备引用。}

VFS 的 \code{ReadAt} 使用显式 offset，不推进共享 FileDescription offset。
否则一个 page fault 会改变进程下一次普通 \code{read} 的位置，把 VM 副作用
泄漏到 Unix I/O 语义。

\topic{文件 VMA 的 split 必须平移 offset}
设 VMA 将三张虚拟页映射到文件 offset 8192。移除中间页后，右段的虚拟起点
增加 8192 字节，其文件起点也必须增加 8192 字节。只修改虚拟半开区间会让
右段重新读取原文件首页。

因此相邻 file VMA 合并的条件比匿名区更严格：

\begin{itemize}
  \item kind、R/W/X 和 map policy 相同；
  \item backing descriptor 与 generation 相同；
  \item 左段文件末端等于右段文件起点；
  \item data length 的尾零视图连续；
  \item owner identifier 相同。
\end{itemize}

中段 unmap 在修改旧描述符前先取得 split descriptor。元数据池耗尽时原图
保持不变，不能留下只有左半次提交的映射。

\topic{ELF 文件大小与内存大小的历史含义}
ELF segment 同时给出 \code{p\_filesz} 和 \code{p\_memsz}。C/C++ 的已初始化
全局对象需要占据文件字节；零初始化对象放入 BSS，只需要在内存中存在，不必
让可执行文件保存大量零。

\begin{center}
\code{memory = file prefix + zero-filled BSS suffix}
\end{center}

fault handler 先清零整页，再只读取 segment 有效 file bytes。普通文件的
3000 字节映射也使用相同尾零原则：前 3000 字节来自文件，后 1096 字节必须
为零。

完整只读页的逻辑内容唯一，可以进入共享 clean cache。部分尾页与可写 private
页使用独立 frame，避免同一个 cache key 同时代表不同的有效长度视图。

\topic{x86-64 \#PF 提供的是原因，不是政策}
页故障 error code 的 P、W/R、U/S、RSVD 与 I/D 位描述硬件观察到的访问：

\begin{itemize}
  \item P=0 才是本阶段可以解析的 not-present；
  \item U/S=1 才允许进入用户 VMA 政策；
  \item RSVD=1 表示页表编码损坏，不能用装页掩盖；
  \item W/R 和 I/D 必须与 VMA writable/executable 一致；
  \item P=1 的写保护错误属于 protection fault，不能当 cache miss。
\end{itemize}

CPU 不知道 inode、ELF BSS 或 \code{MAP\_PRIVATE}。这些政策都由 Kernel 在
读取 CR2、查找 VMA 后决定。成功建立 PTE 后不手动推进 RIP；返回异常现场，
原指令自然重试。

\topic{同一文件页为什么只保留一个 clean cache 项}
当前 cache 初始化时由调用方提供 entry 数组。容量按 managed frames 的
十六分之一缩放，并限制在 256 至 4096 页。fault 热路径不向 Kernel heap
申请任意元数据，资源上界在启动时即可计算。

\begin{pdfdiagram}
  \node[os state] (inactive) {inactive};
  \node[os block,right=of inactive] (clean0) {clean\\refs=0};
  \node[os state,right=of clean0] (cleanN) {clean\\refs>0};
  \draw[os arrow] (inactive) -- node[above,font=\scriptsize]{miss/read} (clean0);
  \draw[os arrow] (clean0) -- node[above,font=\scriptsize]{map} (cleanN);
  \draw[os arrow] (cleanN) to[bend left=22]
      node[below,font=\scriptsize]{unmap} (clean0);
  \draw[os arrow] (clean0) to[bend left=24]
      node[below,font=\scriptsize]{LRU/invalid} (inactive);
\end{pdfdiagram}
\diagramnote{没有 dirty/writeback 状态；零引用 clean 页随时可以从 rootfs
重读，因此可安全丢弃。}

缓存命中增加 mapping reference，PTE 撤销减少 reference。只有 refs=0 的
entry 可被 LRU 选择；全体都被引用时返回明确容量失败。淘汰不能偷走仍由其他
进程页表指向的 frame。

LRU 使用 64 位单调 access generation。每次装入、命中或引用变化都会更新，
容量满时扫描零引用条目并选择最小 generation。O(n) 扫描不是最终高性能结构，
但在最大 4096 entry 的学习内核中让身份、引用和淘汰可以逐项审计。

\topic{MAP\_SHARED 与 MAP\_PRIVATE}
本阶段接受只读 shared 和可写 private：

\begin{osgridlongtable}{L{0.27\textwidth}L{0.25\textwidth}L{0.30\textwidth}}
\textbf{映射} & \textbf{frame} & \textbf{写入语义}\\ \hline
只读 shared & 完整页引用 clean cache & 用户写产生 protection fault\\ \hline
只读 private & 可引用 clean cache & 无写入\\ \hline
可写 private & fault 时读取到独立 frame & 只修改进程副本，不回写\\ \hline
可写 shared & 不支持 & 明确返回错误，不能静默降级\\ \hline
\end{osgridlongtable}

COW 留给下一节。当前 writable private 在第一次 fault 就复制文件
内容，内存效率较低，但不需要提前引入匿名页引用计数和只读 COW PTE 软件位。

\topic{文件修改需要同时处理 PTE 与 cache}
只删除 cache entry 不足以让已建立的 PTE 失效。CPU 命中旧 PTE 时不会重新
进入 Kernel。write/truncate 因此：

\begin{enumerate}[label=\arabic*.]
  \item 扫描仍拥有地址空间的 Process；
  \item 找到相同 FileIdentity 的文件 VMA；
  \item 撤销可重装的只读 PTE，减少 cache 引用；
  \item 更新文件内容或长度；
  \item 失效零引用 cache entry；
  \item 后续访问重新 fault 并观察新内容。
\end{enumerate}

已经退出但尚未 wait 的 Zombie 仍占结果槽，地址空间却已经销毁。扫描条件必须
同时检查 Process active 与 CR3/root physical address 非零。这个跨模块故障
曾使文件写被误报为损坏，只有真实 PID1、Shell、Zombie 与 VM 交错的 QEMU
路径才能稳定暴露。

\topic{日志必须反映事件变化而不是调用次数}
用户的 4096 字节 \code{WriteDescriptor} 会按 256 字节 ABI 上限拆成十六个
系统调用。第一次块写撤销一张 cache page，后十五次没有新 entry。若 logger
只判断累计 invalidation count 仍是二次幂 1，就会打印十六条相同日志。

正确条件是：

\begin{center}
\code{after != before \&\& IsPowerOfTwo(after)}
\end{center}

file fault、cache hit、cache invalidation 与 descriptor block 都只在
1、2、4、8……次变化时采样。QEMU 宿主为每行增加毫秒时间；来宾 PIT 继续
提供独立单调时间。成功探针输出一次语义标记，失败探针输出具名阶段。

\topic{启动布局也是文件装入系统的一部分}
代码增长使 Kernel payload 超过历史 1 MiB 暂存区。Stage 1 的 ATA reader
必须先把容器放到不会覆盖 Kernel segment、临时页表和栈的物理窗口。构建布局
预留 8 MiB \code{0x03600000..0x03E00000} staging，并把 rootfs v2 移到
LBA 32768，即 16 MiB。

汇编、C++ 与 Python 各自需要常量，但不能各自决定事实。Python
\code{boot\_layout.py} 作为镜像工具权威，契约测试扫描 NASM loader、BootInfo
和 Kernel image 工具的边界一致。legacy-fs 的历史 LBA 2048 是另一种格式的
测试区域，不能因数字相同或接近而与生产 rootfs 合并。

\topic{在 Ring 3 连续检查映射语义}
\code{/bin/memory\_probe} 在真实用户态完成：

\begin{enumerate}[label=\arabic*.]
  \item 分块写入一页文件并关闭 fd；
  \item 映射 3000 字节，检查内容和页尾零；
  \item 建立两个 shared 映射并分别触页，要求 cache hit 增长；
  \item 关闭映射所用 fd，继续读取；
  \item 重写文件，要求两个旧 PTE 下次访问都见新内容；
  \item 拒绝 writable shared；
  \item 修改 writable private，再普通读取文件证明没有回写；
  \item unmap 全部区域，核对 file resident、引用与页表回收。
\end{enumerate}

64 MiB、256 MiB 和 64 GiB 使用同一 workload。最低档证明固定元数据不会挤爆
小内存；功能档提供日常完整路径；容量档还建立并销毁 512 个 KernelStack 和
2048 个映射页，最终活动 VMA、对象、描述符与进程均归零。

\topic{测试分层与失败诊断}
单元测试验证唯一 identity、miss/hit、busy invalidation、LRU 和后备
generation；集成测试让两个 AddressSpace 的 PTE 指向同一 frame，并按先后
unmap 证明最后引用所有权；固定种子随机模型执行十万步，每一步对照 active
key、refs、resident 和统计；QEMU 验证 CR2、真实 PTE、ATA/VFS、系统调用
分块与进程交错。

几个具有教学价值的真实故障是：

\begin{itemize}
  \item 新 probe 用旧单次 \code{WriteFile} 写 4096 字节，被 256 字节 ABI
        上限正确拒绝；改用自动分块 \code{WriteDescriptor}；
  \item lazy ELF 后，探针执行自身代码也会新增 resident 页，不能再比较整个
       进程的绝对页数；应比较被测 kind 和至少释放的触页数；
  \item 返回 Ring 3 前仍要求 eager entry PTE，导致
        \code{USER\_RETURN\_REJECTED}；返回验证必须先解析 executable VMA；
  \item rootfs 起点移动后稀疏测试设备仍按旧盘尾，近满测试越界；设备容量改为
       由生产 rootfs start 与 total blocks 推导；
  \item Python ELF 审计仍保留 512 页历史上限；现在与 1 GiB 程序窗口共同
       计算 262143 页。
\end{itemize}

\topic{阶段边界}
当前不包含 dirty page、writeback、writable shared、\code{msync}、
readahead、swap、SMP fault I/O 或动态链接。所有 cache 页都能从 rootfs
重建。下一节才为私有可写页增加共享引用和 COW，并要求 Kernel
\code{CopyToUser} 与用户写 fault 使用同一 break-COW 路径。
这个顺序先回答“相同文件页怎样成为唯一 clean 对象”，再回答“相同私有页
怎样在第一次写时分裂”，避免一次页故障同时调试两种共享关系。
